%% kapa: exam=1 round=33 year=2022 period=2 form=A subject=accounting subject_ko=회계학 seq=36 number=76
%% source: Q-Net 감정평가사 기출문제 (cst00309) artlSeq=5209997 file=raw/1008/5209997/2022년도 제33회 감정평가사 1차시험 2교시(감정평가 관계법규, 회계학).pdf
%% conversion: kordoc 4.24.51 -> tools/gen1.py (auto); review=proofread-sonnet; flags=none
\begin{question}{76}
제품 A와 B를 생산·판매하고 있는 (주)감평의 20×1년 제조간접원가를 활동별로 추적한 자료는 다음과 같다.
\begin{framedtext}
{\small\begin{tabular}{@{}llrrr@{}}
 & 원가동인 & 제품 A & 제품 B & 추적가능원가 \\
자재주문 & 주문횟수 & 20회 & 35회 & ￦55 \\
품질검사 & 검사횟수 & 10회 & 18회 & 84 \\
기계수리 & 기계가동시간 & 80시간 & 100시간 & 180 \\
\end{tabular}\par}
\end{framedtext}
제조간접원가를 활동기준으로 배부하였을 경우 제품 A와 B에 배부될 원가는?
\begin{choices}
  \item 제품 A ￦100 \quad 제품 B ￦219
  \item 제품 A ￦130 \quad 제품 B ￦189
  \item 제품 A ￦150 \quad 제품 B ￦169
  \item 제품 A ￦189 \quad 제품 B ￦130
  \item 제품 A ￦219 \quad 제품 B ￦100
\end{choices}
\end{question}
\answer{2}
